\documentclass[a4paper]{article}
\usepackage[english]{babel}
\usepackage[utf8]{inputenc}
\usepackage{textcomp}
\usepackage{amsmath}
\usepackage{amssymb}
\usepackage{gensymb}
\usepackage{physics}
\usepackage{graphicx}
\usepackage{xcolor}
\usepackage{array}
\usepackage{tabularx}
\usepackage{tikz}
\usepackage{pgfplots}
\usepackage{framed}
\usepackage{xfrac}
\usepackage{epstopdf}
\usepackage[most]{tcolorbox}
\usepackage{fix-cm}
\usepackage{cancel}
\usepackage{pagecolor}
\usepackage[margin=0.5in]{geometry}
\usetikzlibrary{quotes,angles}
\usetikzlibrary{arrows.meta,arrows}
\usetikzlibrary{decorations.pathreplacing}
\usetikzlibrary{calc}
\usetikzlibrary{external}
\tikzexternalize[prefix=tikz/ch20figures/]
\usepgfplotslibrary{fillbetween}

\let\phi\varphi
\let\epsilon\varepsilon
\let\bf\textbf
\let\la\langle
\let\ra\rangle
\pgfplotsset{compat=1.18}
\newcommand\der[2]{\frac{d #1}{d #2}}
\newcommand\Deltat{\Delta t}
\newcommand\rads{\text{ rad\;s}^{-1}}
\newcommand\radss{\text{ rad\;s}^{-2}}
\newcommand\rad{\text{ rad}}
\newcommand\s{\text{ s}}
\newcommand\m{\text{ m}}
\newcommand\km{\text{ km}}
\newcommand\J{\text{ J}}
\newcommand\Nm{\text{ Nm}}
\newcommand\ms{\text{ ms}^{-1}}
\newcommand\mss{\text{ ms}^{-2}}
\newcommand\kg{\text{ kg}}
\newcommand\kgms{\text{ kg\;ms}^{-1}}
\newcommand\kgmm{\text{ kg\;m}^{2}}
\newcommand\kgmms{\text{ kg\;m}^2\text{s}^{-1}}
\newcommand{\ih}{\hat{\imath}}
\newcommand{\jh}{\hat{\jmath}}
\newcommand{\AxisRotator}[1][rotate=0]{%
    \tikz [x=0.25cm,y=0.60cm,line width=.2ex,-stealth,#1] \draw (0,0) arc (-150:150:1 and 1);%
}
\def\centerarc[#1](#2)(#3:#4:#5){\draw[#1] ($(#2)+({#5*cos(#3)},{#5*sin(#3)})$) arc (#3:#4:#5)}
% Syntax: \centerarc[draw options] (center) (initial angle:final angle:radius);

\title{Electric Charge, Force, \& Field}
\author{Essential University Physics Vol. 2, 4e.}
\date{}

\definecolor{fg0}{HTML}{fbf1c7}

\makeatletter % change only the display of \thepage, but not \thepage itself:
\patchcmd{\ps@plain}{\thepage}{\textcolor{fg1}{\thepage}}{}{}
\makeatother

\begin{document}
\pagestyle{plain}

\definecolor{bg0}{HTML}{282828}
\definecolor{bg0_h}{HTML}{1d2021}
\definecolor{bg0_s}{HTML}{32302f}
\definecolor{bg1}{HTML}{3c3836}
\definecolor{bg2}{HTML}{504945}
\definecolor{bg3}{HTML}{665c54}
\definecolor{bg4}{HTML}{7c6f64}
\definecolor{fg1}{HTML}{ebdbb2}
\definecolor{fg2}{HTML}{d5c4a1}
\definecolor{fg3}{HTML}{bdae93}
\definecolor{fg4}{HTML}{a89984}
\definecolor{gbaqua}{HTML}{8ec07c}
\definecolor{gbaqua2}{HTML}{689d6a}
\definecolor{gbred}{HTML}{fb4934}
\definecolor{gbred2}{HTML}{cc241d}
\definecolor{gbgreen}{HTML}{b8bb26}
\definecolor{gbgreen2}{HTML}{98971a}
\definecolor{gbyellow}{HTML}{fabd2f}
\definecolor{gbyellow2}{HTML}{d79921}
\definecolor{gbblue}{HTML}{83a598}
\definecolor{gbblue2}{HTML}{458588}
\definecolor{gbpurple}{HTML}{d3869b}
\definecolor{gbpurple2}{HTML}{b16286}
\definecolor{gborange}{HTML}{fe8019}
\definecolor{gborange2}{HTML}{d65d0e}
\definecolor{gbgray}{HTML}{a89984}
\definecolor{gbgray2}{HTML}{928374}
\color{fg0}
\pagecolor{bg0_s}
\colorlet{shadecolor}{gbaqua2}

\tikzstyle{mybox} = [draw=fg1,fill=bg0_s,very thick,rectangle,rounded corners, inner sep=10pt, inner ysep=10pt]
\tikzstyle{fancytitle} = [draw=fg1,fill=gbaqua2,very thick,text=bg0_s,font=\bfseries,rounded corners]

\setcounter{section}{20}
\maketitle

\subsection{Electric Charge}
Electric charge is an intrinsic property of protons and electrons and comes in two varieties, \textit{positive} and \textit{negative}. These names are useful as the total charge (net charge) of an object is the algebraic sum of its constituent charges. Like charges repel, opposite charges attract, this constitutes a qualitative description of the electric force.
\vspace{1mm}\\
\bf{Quantities of Charge}
\vspace{1mm}\\
All electrons carry the same charge, as do all protons. The charge of a proton is \textit{exactly} the same magnitude as the electron's but with opposite sign. The magnitude of the electron or proton charge is the elementary charge $e$, and is quantized. Elementary particle theories show that the fundamental charge is actually $\frac{1}{3}e$ and resides on quarks, the building blocks of protons and neutrons among other particles. Quarks always join together to form particles with integer multiples of the full elementary charge $e$, and it seems impossible to isolate individual quarks.
\vspace{1mm}\\
The SI unit of charge is the coulomb (C), named after the French physicist Charles Augustin de Coulomb. From the late 19\textsuperscript{th} to early 21\textsuperscript{st} century, the coulomb was defined in terms of electric current and time, a definition that was difficult to implement. The 2019 revision of the SI defined the coulomb more simply by defining the elementary charge as $e = 1.602176634 \times 10^{-19}$ C. The coulomb is therefore the number of elementary charges equal to the inverse of this number: C $\approx 6.24 \times 10^{19}\ e$.
\vspace{1mm}\\
\bf{Charge Conservation}
\vspace{1mm}\\
Electric charge is a conserved quantity, the net charge in a closed region remains constant. Charged particles may be created or annihilated, but always in pairs of equal and opposite charge. The net charge is always the same.
\subsection{Coulomb's Law}
Attraction and repulsion of electric charges implies a force. Joseph Priestley and Charles Augustin de Coulomb investigated this force in the late 1700s and found that the force between two charges acts along the line joining them, with the magnitude proportional to the product of the charges and inversely proportional to the square of the distance between them. Coulomb's law summarizes these results:
\begin{equation}
    \vec{F}_{12} = \frac{kq_1q_2}{r^2}\hat{r}
\end{equation}
where $\vec{F}_{12}$ is the force charge $q_1$ exerts on $q_2$ and $r$ is the distance between the charges. In SI, the proportionality constant $k$ has the approximate value $9 \times 10^9\; \text{N}\; \text{m}^2\; \text{C}^{-2}$. Force is a vector, and $\hat{r}$ is a unit vector that helps determine its direction, pointing \textit{from} $q_1$ \textit{toward} $q_2$. Reversing the roles of $q_1$ and $q_2$, $\vec{F}_{21}$ has the same magnitude as $\vec{F}_{12}$ but the opposite direction, thus Coulomb's law obeys Newton's third law. The force is in the same direction as the unit vector when the charges have the same sign, and opposite the unit vector when the charges have opposite signs, accounting for the fact that like charges repel and opposite charges attract.
\begin{center}
    \begin{tikzpicture}
        \draw[->,very thick,-latex,gbblue2] (1,0.9)--node[below,color=fg1]{$\vec{F}_{12}$} (3,0.9);
        \draw[->,very thick,-latex,gbaqua2] (1,1.1)--node[above,color=fg1]{$\hat{r}$}(2,1.1);
        \filldraw[gbblue] (-1,1) circle (0.15);
        \filldraw[gbblue] (1,1) circle (0.15);
        \draw[thin] (-0.9,1)--(-1.1,1);
        \draw[thin] (-1,1.1)--(-1,0.9);
        \draw[thin] (0.9,1)--(1.1,1);
        \draw[thin] (1,1.1)--(1,0.9);
        \draw[thick] (-1,1) circle (0.15);
        \draw[thick] (1,1) circle (0.15);
        \draw[thick] (-1,0.75)--(-1,0);
        \draw[thick] (1,0.75)--(1,0);
        \draw[thick] (-1,0.1)--node[fill=bg0_s]{$r$} (1,0.1);
        \node at (0,-0.5) {(a)};
    \end{tikzpicture}
    \hspace{15mm}
    \begin{tikzpicture}
        \draw[->,very thick,-latex,gbblue2] (1,1)--node[above,color=fg1]{$\vec{F}_{12}$} (-0.85,1);
        \draw[->,very thick,-latex,gbaqua2] (1,1)--node[above,color=fg1]{$\hat{r}$}(2,1);
        \filldraw[gbblue] (-1,1) circle (0.15);
        \filldraw[gbred] (1,1) circle (0.15);
        \draw[thin] (-0.9,1)--(-1.1,1);
        \draw[thin] (-1,1.1)--(-1,0.9);
        \draw[thin] (0.9,1)--(1.1,1);
        \draw[thick] (-1,1) circle (0.15);
        \draw[thick] (1,1) circle (0.15);
        \draw[thick] (-1,0.75)--(-1,0);
        \draw[thick] (1,0.75)--(1,0);
        \draw[thick] (-1,0.1)--node[fill=bg0_s]{$r$} (1,0.1);
        \node at (0,-0.5) {(b)};
    \end{tikzpicture}
\end{center}
\bf{Point Charges \& the Superposition Principle}
\vspace{1mm}\\
Coulomb's law is only strictly true for point charges (charged objects of negligible size) which electrons and protons can be usually be treated as. Any two charged objects can also be considered as such if the distance between them is large compared to their size. Often the electric effects of charge distributions - arrangements of charge spread over space - are more interesting. To find the electric effect of such charge distributions, the effects of two or more charges must be combined.
\begin{center}
    \begin{tikzpicture}[scale=1.5]
        \draw[->,ultra thick,-latex,gbblue2] (0,0)--node[below,color=fg1]{$\vec{F}_{13}$}(1,0);
        \draw[->,ultra thick,-latex,gbblue2] (0,0)--node[right,color=fg1]{$\vec{F}_{\text{net}} = \vec{F}_{13} + \vec{F}_{23}$}({cos(60)},{sin(60)});
        \draw[->,ultra thick,-latex,gbblue2] (0,0)--node[left,color=fg1]{$\vec{F}_{23}$}({-cos(60)},{sin(60)});
        \filldraw[gbblue] (0,0) circle (0.1);
        \draw[thick] (0,0) circle (0.1);
        \node[below] at (0,-0.1){$q_3$};
        \draw[thin] (0,-0.075)--(0,0.075);
        \draw[thin] (-0.075,0)--(0.075,0);
        \filldraw[gbred] ({-1.5*cos(60)},{1.5*sin(60)}) circle (0.1);
        \draw[thin] ({-1.5*cos(60) - 0.075},{1.5*sin(60)})--({-1.5*cos(60) + 0.075},{1.5*sin(60)});
        \draw[thick] ({-1.5*cos(60)},{1.5*sin(60)}) circle (0.1);
        \node[below] at ({-1.5*cos(60)},{1.5*sin(60) - 0.1}){$q_2$};
        \filldraw[gbblue] (-1,0) circle (0.1);
        \draw[thin] (-1.075,0)--(-0.925,0);
        \draw[thin] (-1,0.075)--(-1,-0.075);
        \node[below] at (-1,-0.1){$q_1$};
        \draw[thick] (-1,0) circle (0.1);
    \end{tikzpicture}
\end{center}
The figure shows two charges $q_1$ and $q_2$ constituting a simple charge distribution. To calculate the net force these exert on a third charge $q_3$, calculate the forces $\vec{F}_{13}$ and $\vec{F}_{23}$ and add the two vectors. The force $q_1$ exerts on $q_3$ is unaffected by the presence pf $q_2$ and vice versa. Coulomb's law can be applied separately to the pairs $q_1q_3$ and $q_2q_3$ and the results combined. This fact - that electric forces add vectorially - is called the superposition principle.
\subsection{The Electric Field}
The gravitational field at a point is defined as the gravitational force per unit mass that an object at that point would experience. In this context, $\vec{g}$ can be thought of as the force per unit mass that any object would experience due to gravity. The gravitational field can be pictured as a continuous set of vectors that give the magnitude and direction of the gravitational force per unit mass at each point. The same can be done with the electric force, defining the electric field as the force per unit charge:
\begin{equation}
    \vec{E} = \frac{\vec{F}}{q}
\end{equation}
This equation can be used as a prescription for measuring electric fields. Place a point charge at some location, measure the electric force it experiencesm and divide by the change to get the field.
\begin{center}
    \begin{tikzpicture}[scale=1.5]
        \draw[thick] (0,0)--(2,0);
        \node[right] at (2,0.8){$\vec{g}$};
        \node at (0,-1.125){};
        \foreach \i in {0.1,0.3,...,1.9}{
            \draw[->,very thick,-latex,gbblue2] (\i,0.75)--(\i,0.1);
            \draw[->,very thick,-latex,gbblue2] (\i,1.5)--(\i,0.85);
            \filldraw (\i,0.75) circle (0.05);
            \filldraw (\i,1.5) circle (0.05);
        }
    \end{tikzpicture}
    \hspace{15mm}
    \begin{tikzpicture}[scale=1.5]
        \filldraw[gbblue] (0,0) circle (0.15);
        \draw[thick] (0,0) circle (0.15);
        \draw[thin] (-0.1,0)--(0.1,0);
        \draw[thin] (0,-0.1)--(0,0.1);
        \foreach \i in {22.5,67.5,112.5,...,337.5}{
            \draw[->,very thick,-latex] ({0.5*cos(\i)},{0.5*sin(\i)})--({1.5*cos(\i)},{1.5*sin(\i)});
            \filldraw ({0.5*cos(\i)},{0.5*sin(\i)}) circle (0.05);
            \draw[->,very thick,-latex] ({1.6*cos(\i)},{1.6*sin(\i)})--({2.25*cos(\i)},{2.25*sin(\i)});
            \filldraw ({1.6*cos(\i)},{1.6*sin(\i)}) circle (0.05);
        };
        \foreach \j in {0,45,90,...,315}{
            \draw[->,very thick,-latex] ({0.35*cos(\j)},{0.35*sin(\j)})--({2.25*cos(\j)},{2.25*sin(\j)});
            \filldraw ({0.35*cos(\j)},{0.35*sin(\j)}) circle (0.05);
        }
        \node at ({1.2*cos(58)},{1.2*sin(58)}){$\vec{E}_1$};
        \node at ({1.9*cos(32)},{1.9*sin(32)}){$\vec{E}_2$};
    \end{tikzpicture}
\end{center}
At the point at the tail of $\vec{E}_1$, the electric field is described by the vector $\vec{E}_1$ meaning a point charge $q$ placed there would experience an electric force $q\vec{E}_1$. Farther from the charge at the tail of $\vec{E}_2$, a point charge $q$ would experience a weaker force $q\vec{E}_2$.
\vspace{1mm}\\
If the electric field $\vec{E}$ at a point is known, equation 2 can be rearranged to find the force on any point charge $q$ placed at that point:
\begin{equation*}
    \vec{F} = q\vec{E}
\end{equation*}
If the charge $q$ is positive, the force is in the same direction as the field, if $q$ is negative, then the force is opposite to the direction of the field. The units of electric field are newtons per coulomb (N$\;$C$^{-1}$). Fields of hundreds or thousands of N$\;$C$^{-1}$ are common, while fields of 3 MN$\;$C$^{-1}$ will tear electrons from air molecules.
\vspace{1mm}\\
\bf{The Field of a Point Charge}
\vspace{1mm}\\
Once the field of a charge distribution is known, its effect on other charges can be calculated. Coulomb's law gives the force on a test charge $q_{\text{test}}$ located a distance $r$ from a point charge $q$: $\vec{F} = (\sfrac{kqq_{\text{test}}}{r^2})\hat{r}$, where $\hat{r}$ is a unit vector pointing away from $q$. The electric field arising from $q$ is the force per unit charge, or:
\begin{equation}
    \vec{E} = \frac{\vec{F}}{q_{\text{test}}} = \frac{kq}{r^2}\hat{r}
\end{equation}
Because of its similarity, this equation is also referred to as Coulomb's law. The equation contains no reference to the test charge $q_{\text{test}}$ because the field of $q$ exists independently of any other charge. Since $\hat{r}$ always points away from $q$, the direction of $\vec{E}$ is radially outward if $q$ is positive and radially inward if $q$ is negative.
\begin{center}
    \begin{tikzpicture}
        \filldraw[gbred] (0,0) circle (0.1);
        \draw[thick] (0,0) circle (0.1);
        \draw[thin] (-0.06,0)--(0.06,0);
        \foreach \i in {0,45,90,...,315}{
            \draw[->,very thick,-latex] ({cos(\i)},{sin(\i)})--({0.2*cos(\i)},{0.2*sin(\i)});
            \draw[->,very thick,-latex] ({1.75*cos(\i)},{1.75*sin(\i)})--({1.1*cos(\i)},{1.1*sin(\i)});
            \filldraw ({cos(\i)},{sin(\i)}) circle (0.05);
            \filldraw ({1.75*cos(\i)},{1.75*sin(\i)}) circle (0.05);
        }
        \foreach \j in {22.5,67.5,112.5,...,337.5}{
            \draw[->,very thick,-latex] ({1.6*cos(\j)},{1.6*sin(\j)})--({0.8*cos(\j)},{0.8*sin(\j)});
            \filldraw ({1.6*cos(\j)},{1.6*sin(\j)}) circle (0.05);
        }
        \node at (0,-2.15){Field vectors for a negative point charge};
    \end{tikzpicture}
\end{center}
\newpage
\subsection{Fields of Charge Distribution}
Since the electric force obeys the superposition principle, so does the electric field, meaning that the field of a charge distribution is the vector sum of the fields of the individual point charges making up the distribution:
\begin{equation}
    \vec{E} = \vec{E}_1 + \vec{E}_2 + \vec{E}_3 + ... = \sum_i \vec{E}_i = \sum_i \frac{kq_i}{r_i^2}\hat{r}_i
\end{equation}
In this equation, the $\vec{E}_i$'s are the electric fields of the point charges $q_i$ located at distance $r_i$ from the field point (point where the field is being evaluated). The unit vectors $\hat{r}_i$ point from each point charge toward the field point. In theory this equation gives the electric field of any charge distribution, although this is often complicated unless the distribution contains relatively few points.
\begin{center}
    \begin{tikzpicture}
        \node[mybox,text=fg1](box){
            \begin{minipage}{0.95\textwidth}
                Two protons are 3.6 nm apart. Find the electric field at a point between them, 1.2 nm from one of the protons. Then find the force on an electron at this point.
                \begin{center}
                    \begin{tikzpicture}
                        \draw[<->,latex-latex] (-0.5,0)--(4.1,0) node[right]{$x$};
                        \draw[->,thick,-latex,gbblue2] (0,0)--node[above,yshift=-2mm,text=fg1]{$\hat{r}_1$} (0.7,0);
                        \draw[->,thick,-latex,gbblue2] (3.6,0)--node[above,yshift=-2mm,text=fg1]{$\hat{r}_2$} (2.9,0);

                        \draw[<->,latex-latex] (0,-0.25)--(3.6,-0.25);
                        \draw[thick,bg0_s] (1.2,-0.25)--(2.4,-0.25);
                        \node at (1.8,-0.25) {\small 3.6 nm};

                        \draw[<->,latex-latex] (0,-0.5)--node[below,yshift=2mm]{\small 1.2 nm}(1.2,-0.5);
                        \draw[<->,latex-latex] (1.2,-0.5)--node[below,yshift=2mm]{\small 2.4 nm}(3.6,-0.5);

                        \filldraw[gbblue2] (0,0) circle (0.15);
                        \draw[thick] (0,0) circle (0.15);
                        \draw (-0.09,0)--(0.09,0);
                        \draw (0,-0.09)--(0,0.09);

                        \filldraw[gbblue2] (3.6,0) circle (0.15);
                        \draw[thick] (3.6,0) circle (0.15);
                        \draw (3.51,0)--(3.69,0);
                        \draw (3.6,-0.09)--(3.6,0.09);

                        \filldraw (1.2,0) circle (0.04) node[above,yshift=-2mm]{$P$};
                    \end{tikzpicture}
                \end{center}
                The field point $P$ is identified as being 1.2 nm from one proton. By letting the line between the protons define the $x$-axis, the unit vector $\hat{r}_1$ becomes $\ih$, and $\hat{r}_2$ becomes $-\ih$. 
                \vspace{1mm}\\
                The charge $q$ of both protons is the elementary charge $e$ and the charge of an electron is $-e$, evaluating using equation 4 gives
                \begin{align*}
                    \vec{E} &= \vec{E}_1 + \vec{E}_2 = \frac{k_e e}{r_1^2}\ih + \frac{k_e e}{r_2^2}(-\ih) = k_e e\left(\frac{1}{r_1^2} - \frac{1}{r_2^2}\right)\ih \\
                    &= \left(8.99 \cdot 10^9\right)\left(1.6\cdot 10^{-19}\right)\left(\frac{1}{1.2^2} - \frac{1}{2.4^2}\right)\ih = 750\;\ih \text{ MN}\;\text{C}^{-1}\\
                    \vec{F} &= q\vec{E} = -e\vec{E} = \left(-1.6\cdot 10^{-19}\right)\left(7.5\cdot 10^{8}\right) = -1.2\cdot 10^{-10}\ih \text{ N}
                \end{align*}
                Therefore, the force on an electron at point $P$ is 0.12 nN in the $-\ih$ direction.
            \end{minipage}
        };
        \node[fancytitle,right=10pt] at (box.north west) {Example: Field From Two Protons};
    \end{tikzpicture}
\end{center}

\newpage
\noindent\bf{The Electric Dipole}
\vspace{1mm}\\
The electric dipole is one of the most important charge distribution and consists of two point charges of equal magnitude and opposite sign. Many molecules have dipoles, which makes them essential in explaining molecular behavior.
\begin{center}
    \includegraphics{tikz/ch20figures/H2O-dipole.eps}
\end{center}
The net charge of a water molecule is 0 however, due to the different electronegativities of oxygen and hydrogen, a dipole is formed from partial positive/negative charges.

\begin{center}
    \begin{tikzpicture}
        \node[mybox,text=fg1] (box) {
            \begin{minipage}{0.95\textwidth}
                A molecule may be approximately modeled as a positive charge $q$ at $x = a$ and a negative charge $-a$ at $x  = -a$. Evaluate the electric field on the $y$-axis, and find an approximate expression valid at large distances ($|y| \gg a$).
                \begin{center}
                    \begin{tikzpicture}[scale=1.75]
                        \draw[<->,latex-latex] (-1,0)--(1,0) node[right,xshift=-2mm]{$x$};
                        \draw[->,-latex] (0,-0.25)--(0,1.5) node[above,yshift=-2mm]{$y$};
                        \draw[dashed,thin,fg3] (-0.5,0)--(0,0.9)--(0.5,0);
                        
                        \filldraw[gbred] (-0.5,0) circle (0.1);
                        \draw (-0.5,0) circle (0.1);
                        \draw[thin] (-0.55,0)--(-0.45,0);

                        \filldraw[gbblue2] (0.5,0) circle (0.1);
                        \draw (0.5,0) circle (0.1);
                        \draw[thin] (0.55,0)--(0.45,0);
                        \draw[thin] (0.5,0.05)--(0.5,-0.05);

                        \draw[->,thick,-latex,gbblue2] (0,0.9)--({-0.5*cos(60.945)},{0.9 + 0.5*sin(60.945)}) node[left,yshift=2mm,xshift=5mm,text=fg1]{$\vec{E}_{+}$};
                        \draw[->,thick,-latex,gbblue2] (0,0.9)--({-0.5*cos(60.945)},{0.9 - 0.5*sin(60.945)}) node[left,yshift=4.5mm,xshift=5mm,text=fg1]{$\vec{E}_{-}$};
                        \draw[->,thick,-latex,gbblue2] (0,0.9)--(-0.5,0.9) node[left,text=fg1,xshift=2mm]{$\vec{E}$};
                        \filldraw (0,0.9) circle (0.025);
                        \node at (-0.5,0.45) {$r$};
                        \node at (0.5,0.45) {$r$};
                        \node at (-0.7,0.2) {$-q$};
                        \node at (0.7,0.2) {$q$};
                    \end{tikzpicture}
                \end{center}
                The individual unit vectors point from the two charges toward the field point. The negative charge contributes a field opposite its unit vector; individual fields are indicated. Symmetry makes the $y$-components cancel, giving a field in the $-x$ direction, so only the $x$-components of the unit vectors are needed.
                \vspace{1mm}\\
                These are $\hat{r}_{x^-} = \sfrac{a}{r}$ for the negative charge at $x = -a$ and $\hat{r}_{x^+} = -\sfrac{a}{r}$ for the positive charge at $x = a$. \\
                When evaluated using equation 4:
                \begin{align*}
                    \vec{E} = \frac{k_e(-q)}{r^2}\left(\frac{a}{r}\right)\ih + \frac{k_eq}{r^2}\left(-\frac{a}{r}\right)\ih = -\frac{2k_eqa}{(a^2 + y^2)^{\sfrac{3}{2}}}
                \end{align*}
                In the last step the substitution $r = \sqrt{a^2 + y^2}$ is used. For $|y| \gg a$, $a^2$ can be neglected compared with $y^2$, giving
                \begin{align*}
                    \vec{E} = -\frac{2k_eqa}{|y|^3}\ih \hspace{5mm} (|y| \gg a)
                \end{align*}
            \end{minipage}
        };
        \node[fancytitle,right=10pt] at (box.north west) {Example: Electric Dipole - Modeling a Molecule};
    \end{tikzpicture}
\end{center}
The dipole field at large distances decreases as the inverse cube of distance $(\vec{E} \propto \sfrac{1}{r^3})$. This is because the net charge of the dipole is 0, its field comes entirely from the separation of opposite charges. Because of this, the dipole field is not exactly zero, but it weaker and more localized than that of a point charge. At large distances, a dipole's electric properties are entirely characterized by its dipole moment $p$, which is defined as the product of its charge $q$ and the distance $d$ between the charges.
\begin{equation}
    p = qd
\end{equation}
In the above example, $d = 2a$, so the dipole moment was $p = 2qa$. The field can be written in terms of the dipole moment:
\begin{align}
    \begin{split}
        \vec{E} = -\frac{k_ep}{|y|^3}\ih \\
        \vec{E} = \frac{2k_ep}{|x|^3}\ih 
    \end{split}
\end{align}
These equations show that the field along the axis of the dipole at a given distance is twice as strong as along the bisector. 
\vspace{1mm} \\

\noindent\bf{Continuous Charge Distributions}
\vspace{1mm}\\
If the charge distribution extends throughout a volume, it is described in terms of volume charge density $\rho$ with units C$\;$m$^{-3}$. For charge distributions spread over lines or surfaces, surface charge density $\sigma$ (C$\;$m$^{-2}$), and linear charge density $\lambda$ (C$\;$m$^{-1}$) are used. To calculate the field of a continuous distribution, the charged region is divided into infinitesimal charge elements $dq$. Each $dq$ produces a field $d\vec{E}$ ($d\vec{E} = \frac{k_edq}{r^2}\hat{r}$), the sum of all the $d\vec{E}$'s is the integral
\begin{equation}
    \vec{E} = \int d\vec{E} = \int \frac{k_edq}{r^2}\hat{r}
\end{equation}

\newpage
\begin{center}
    \begin{tikzpicture}
        \node[mybox,text=fg1] (box) {
            \begin{minipage}{0.95\textwidth}
                A ring of radius $a$ has a charge $Q$ distributed evenly over the ring. Find an expression for the electric field at any point on the axis of the ring.
                \vspace{1mm}\\
                Take the $x$-axis to coincide with the ring axis with the center of the ring at $x = 0$. The $y$-components of the field contributions from pairs of charge elements on opposite sides of the ring cancel so the net field points in the $+x$ direction (for $x > 0$), and only the $x$-components of the unit vectors are needed, namely, $\hat{r}_x = \cos(\theta) = \frac{x}{r}$
                \begin{center}
                    \begin{tikzpicture}[scale=1.6]
                        \filldraw[gbred] (0.029,0.9)--(0.029,1.1)--(-0.029,1.1)--(-0.029,0.9);
                        \filldraw[gbred] (0.029,1.1)--(0.029,0.9)--(-0.029,0.9)--(-0.029,1.1);
                        \filldraw[gbred] (0.029,-0.9)--(0.029,-1.1)--(-0.029,-1.1)--(-0.029,-0.9);
                        \filldraw[gbred] (0.029,-1.1)--(0.029,-0.9)--(-0.029,-0.9)--(-0.029,-1.1);
                        \draw[densely dashed,fg3] (0,0)--(0,-0.33);
                        \draw[densely dashed,fg3] (2,0)--(2,-0.33);
                        \draw[thick] (0,0) ellipse (0.55 and 1.1);
                        \draw[thick] (0,0) ellipse (0.4 and 0.9);
                        \filldraw (0,0) circle (0.04) coordinate (o);
                        \draw (-0.03,1.1)--(-0.03,0.9);
                        \draw (0.03,1.1)--(0.03,0.9);
                        \draw (-0.03,-1.1)--(-0.03,-0.9);
                        \draw (0.03,-1.1)--(0.03,-0.9);
                        \draw (0,1) coordinate (a) node[above]{$dq$};
                        \draw (0,-1) coordinate (b) node[above]{$dq$};
                        \draw[<->,latex-latex] (0,0.04)--node[left,xshift=2mm]{$a$}(0,0.9);
                        \draw (-0.6,0)--(-0.55,0);
                        \draw (-0.4,0)--(o);

                        \draw (2,0) coordinate (P);

                        \draw ({2 + 0.75*cos(26.565)},{0.75*sin(-26.565)}) coordinate (c);
                        \draw ({2 + 0.75*cos(26.565)},{0.75*sin(-26.565)}) coordinate (c);
                        \draw (3,0) coordinate (d);
                        \draw pic["$\theta$",draw=fg1,-,thick,angle eccentricity=1.5,angle radius=6mm]{angle=c--P--d};
                        \draw[<->,latex-latex] (0,-0.25)--node[above,yshift=-2.7mm]{$x$}(2,-0.25);

                        \draw[->,-latex,thick,gbblue2] (2,0)--({2 + 0.75*cos(26.565)},{0.75*sin(26.565)});
                        \draw[->,-latex,thick,gbblue2] (2,0)--({2 + 0.75*cos(26.565)},{0.75*sin(-26.565)});
                        \draw[->,-latex,thick,gbblue2] (2,0)--(3,0) node[right,xshift=-2.5mm,text=fg1]{$dE_x$};
                        
                        \draw (o)--(P);
                        \draw[dashed,fg3] (a)--(P)--(b);
                        \filldraw (P) circle (0.04) node[above,yshift=-2mm]{$P$};

                        \draw pic["$\theta$",draw=fg1,-,thick,angle eccentricity=1.5,angle radius=6mm]{angle=a--P--o};

                        \draw[->,-latex,thick,gbblue2] (a)--node[above,text=fg1,yshift=-1mm]{$\hat{r}$}({0.75*cos(26.565)},{1 - 0.75*sin(26.565)});
                        \draw[->,-latex,thick,gbblue2] (0,-1)--node[below,text=fg1,yshift=1mm]{$\hat{r}$}({0.75*cos(26.565)},{0.75*sin(26.565) - 1});
                        
                        \node[rotate=-26.565] at (1.25,0.6) {$r = \sqrt{x^2 + a^2}$};
                        \node[below,yshift=3mm] at (c) {$d\vec{E}$};
                        \node[above,yshift=-2.5mm] at ({2 + 0.75*cos(26.565)},{0.75*sin(26.565)}) {$d\vec{E}$};
                    \end{tikzpicture}
                \end{center}
                Each charge element $dq$ contributes the same amount to the field, $\displaystyle dE_x = \frac{k_edq}{r^2}\hat{r}_x = \frac{k_exdq}{r^3}$\\
                Writing $\displaystyle r = \sqrt{x^2 + a^2}$ as $(x^2 + a^2)^{\sfrac{1}{2}}$, the integral becomes
                \begin{align*}
                    E = \int\limits_{\text{ring}} dE_x = \int\limits_{\text{ring}} \frac{k_ex\; dq}{(x^2 + a^2)^{\sfrac{3}{2}}} = \frac{k_ex}{(x^2 + a^2)^{\sfrac{3}{2}}} \int\limits_{\text{ring}} dq
                \end{align*}
                Because the field point $P$ is fixed, its coordinate $x$ is a constant in the integration. The remaining integral is the sum of all the charge elements on the ring, the total charge $Q$.
                \begin{align*}
                    E = \frac{k_eQx}{(x^2 + a^2)^{\sfrac{3}{2}}}
                \end{align*}
            \end{minipage}
        };
        \node[fancytitle,right=10pt] at (box.north west) {Example: Evaluating The Field - A Charged Ring};
    \end{tikzpicture}
\end{center}

\begin{center}
    \begin{tikzpicture}
        \node[mybox,text=fg1](box){
            \begin{minipage}{0.95\textwidth}
                A long, straight electric power line coincides with the $x$-axis and carries a uniform line charge density $\lambda$ (unit = C$\;$m$^{-1}$). Find the electric field on the $y$-axis using the approximation that the wire is infinitely long. \\
                The field point $P$ is placed along the $y$-axis and the wire is divided into infinitesimal charge elements $dq$. The $x$-components of the fields $d\vec{E}$ contributed by elements $dq$ on opposite sides of the $y$-axis cancel, so only the $y$-components of each unit vector are needed. $\hat{r}_y = \sin(\theta) = \sfrac{y}{r}$ where $r = \sqrt{x^2 + y^2}$
                \begin{center}
                    \begin{tikzpicture}[scale=1.5]
                        \draw[<->,latex-latex] (0,-0.25)--node[below,yshift=2mm]{$x$}(1,-0.25);
                        
                        \draw (0,0.75) coordinate (P);
                        \draw[->,-latex,thick,gbblue2] (P)--({0.5*cos(143.13)},{0.75 + 0.5*sin(143.13)}) node[left,text=fg1,xshift=2.5mm]{$d\vec{E}$};
                        \draw[->,-latex,thick,gbblue2] (P)--({-0.5*cos(143.13)},{0.75 + 0.5*sin(143.13)}) node[right,text=fg1,xshift=-2.5mm]{$d\vec{E}$};
                        \draw[dashed,fg3] (1,0)--(P)--(-1,0);
                        \draw[->,-latex,thick,gbblue2] (P)--(0,1.35) node[above,yshift=-3mm,text=fg1]{$dE_y$};
                        \filldraw (P) circle (0.05) node[left,xshift=1mm]{$P$};

                        \draw[->,-latex,thick,gbblue2] (1,0)--({1 + 0.5*cos(143.13)},{0.5*sin(143.13)}) node[above,yshift=-3mm,xshift=2mm,text=fg1]{$\hat{r}$};;
                        \draw[->,-latex,thick,gbblue2] (-1,0)--({-0.5*cos(143.13) - 1},{0.5*sin(143.13)}) node[above,yshift=-3mm,xshift=-2mm,text=fg1]{$\hat{r}$};
                        \node at (1,0.7){$r = \sqrt{x^2 + y^2}$};

                        \draw[ultra thick,gbred] (1.025,-0.1)--(1.025,0) node[above,text=fg1,yshift=-3mm,xshift=1.5mm]{$dq$};
                        \draw[thin] (1,0)--(1,-0.1);
                        \draw[thin] (1.05,0)--(1.05,-0.1);
                        \draw[ultra thick,gbred] (-1.025,-0.1)--(-1.025,0) node[above,text=fg1,yshift=-3mm,xshift=-1.5mm]{$dq$};
                        \draw[thin] (-1,0)--(-1,-0.1);
                        \draw[thin] (-1.05,0)--(-1.05,-0.1);

                        \draw (-1.5,0)--(1.5,0);
                        \draw (-1.5,-0.1)--(1.5,-0.1);
                        \draw (0,-0.5)--(0,0.75);
                        \draw[<->,latex-latex] (-0.1,0)--node[left,xshift=2mm,yshift=-1mm]{$y$}(-0.1,0.7);
                    \end{tikzpicture}
                \end{center}
                The wire has charge density $\lambda$ C$\;$m$^{-1}$, so if a charge element has length $dx$, its charge is $dq = \lambda dx$
                \begin{align*}
                    dE_y = \frac{k_edq}{r^2}\hat{r} = \frac{k_e\lambda dx}{r^2} \frac{y}{r} = \frac{k\lambda y}{(x^2 + y^2)^{\sfrac{3}{2}}}dx
                \end{align*}
                Where $r = \sqrt{x^2 + y^2}$. Since the $x$-components cancel, the $y$-components can be integrated to get the net field
                \begin{align*}
                    E = E_y &= \int_{-\infty}^{\infty} \frac{k_e\lambda y}{(x^2 + y^2)^{\sfrac{3}{2}}}dx = k_e\lambda y \int_{-\infty}^{\infty} \frac{1}{(x^2 + y^2)^{\sfrac{3}{2}}}dx \\
                    &= k_e\lambda y \left[\frac{x}{y^2\sqrt{x^2 + y^2}}\right]_{-\infty}^{\infty} = k_e\lambda y \left[\frac{1}{y^2}-\left(\frac{1}{y^2}\right)\right] = \frac{2k_e\lambda}{y}
                \end{align*}
                The result is the field's magnitude in the radial direction away from the wire for $+\lambda$ and toward the wire for $-\lambda$
            \end{minipage}
        };
        \node[fancytitle,right=10pt] at (box.north west){Example: Line Charge - A Power Line's Field};
    \end{tikzpicture}
\end{center}

\subsection{Matter in Electric Fields}
The definition of the electric field governs the motion of a single charge in an electric field. Combining the equation $\vec{F} = q\vec{E}$ and Newton's law $\vec{F} = m\vec{a}$ gives the acceleration of a particle with charge $q$ and mass $m$ in an electric field $\vec{E}$.
\begin{equation}
    \vec{a} = \frac{q\vec{E}}{m}
\end{equation}
The charge to mass ratio $\frac{q}{m}$ of a particle determines its response to an electric field. When the electric field is uniform, problems involving the motion of charged particles reduce to constant acceleration problems from physics I.
\begin{center}
    \begin{tikzpicture}
        \node[mybox,text=fg1](box){
            \begin{minipage}{0.95\textwidth}    
                Two oppositely charged, curved metal plates establish an electric field given by $E = E_0\frac{b}{r}$ where $E_0$ and $b$ are constants representing the electric field and length respectively. The field points toward the center of curvature, and $r$ is the distance from the center.
                \vspace{1mm} \\
                Find an expression for the speed $v$ with which a proton entering vertically from below will leave the device moving horizontally. \\
                Here, we want uniform circular motion so equation 8 is written with the given field and $\displaystyle a = \frac{v^2}{r}$.
                \begin{align*}
                    a = \frac{v^2}{r} = \frac{eE}{m} = \frac{eE_0b}{mr}
                \end{align*}
                Solving for $v$ gives:
                \begin{align*}
                    v^2 = \frac{eE_0b}{m} \hspace{10mm} v = \sqrt{\frac{eE_0b}{m}}
                \end{align*}
            \end{minipage}
        };
        \node[fancytitle,right=10pt] at (box.north west) {Example: Particle Motion - Electrostatic Analyzer};
    \end{tikzpicture}
\end{center}
\bf{Dipoles in Electric Fields}

\end{document}